% !TeX program = tectonic
\documentclass[11pt]{article}
\IfFileExists{alius-guidelines-style.tex}{\usepackage[a4paper,top=27mm,bottom=24mm,left=25mm,right=25mm]{geometry}
\usepackage[T1]{fontenc}
\usepackage[utf8]{inputenc}
\usepackage[scaled=0.96]{helvet}
\renewcommand{\familydefault}{\sfdefault}
\usepackage{array}
\usepackage{enumitem}
\usepackage{fancyhdr}
\usepackage[hidelinks]{hyperref}
\usepackage{microtype}
\usepackage{needspace}
\usepackage{tabularx}
\usepackage{xcolor}
\usepackage{xurl}

\definecolor{aliusgreen}{HTML}{13823A}
\definecolor{aliusgray}{HTML}{7A7A7A}
\definecolor{aliusline}{HTML}{9A9A9A}

\hypersetup{
  colorlinks=true,
  linkcolor=aliusgreen,
  urlcolor=aliusgreen
}

\setlength{\parindent}{0pt}
\setlength{\parskip}{0.72em}
\setlength{\tabcolsep}{0pt}
\raggedbottom
\sloppy
\urlstyle{same}

\newcommand{\ALIUSFooterTitle}{ALIUS Bulletin Guidelines (2022)}

\pagestyle{fancy}
\fancyhf{}
\fancyfoot[L]{\footnotesize\color{aliusgray}\ALIUSFooterTitle}
\fancyfoot[R]{\footnotesize\color{aliusgray}\thepage}
\renewcommand{\headrulewidth}{0pt}
\renewcommand{\footrulewidth}{0.45pt}

\setlist[itemize]{
  leftmargin=1.35em,
  itemsep=0.18em,
  topsep=0.15em,
  parsep=0em
}
\setlist[enumerate]{
  leftmargin=2.05em,
  itemsep=0.42em,
  topsep=0.2em,
  parsep=0em
}
\setlist[description]{
  leftmargin=0em,
  labelsep=0.55em,
  itemsep=0.35em,
  topsep=0.2em,
  font=\normalfont\bfseries
}

\newcommand{\guideDivider}{%
  \par\noindent{\color{aliusline}\rule{\textwidth}{0.45pt}}\par
}

\newcommand{\guideMeta}[2]{%
  {\footnotesize\bfseries\MakeUppercase{#1}}\par
  {\footnotesize #2}\par\vspace{0.65em}
}

\newcommand{\guideHeader}[5]{%
  \thispagestyle{fancy}%
  \vspace*{1mm}
  \begin{center}
    {\Large\bfseries\MakeUppercase{ALIUS Bulletin}\par}
    \vspace{0.25em}
    {\normalsize *\par}
    \vspace{0.25em}
    {\normalsize\bfseries\MakeUppercase{Guidelines For}\par}
    \vspace{0.15em}
    {\LARGE\bfseries\MakeUppercase{#1}\par}
  \end{center}
  \vspace{1.05em}
  \begin{tabularx}{\textwidth}{@{}>{\raggedright\arraybackslash}p{0.55\textwidth}!{\hspace{1.35em}{\color{aliusline}\vrule width 0.5pt}\hspace{1.35em}}>{\raggedright\arraybackslash}X@{}}
    {\large\itshape #2} &
    \guideMeta{Format}{#3}
    \guideMeta{Editorial Route}{#4}
    \guideMeta{Contact}{#5}
  \end{tabularx}
  \vspace{0.75em}
  \guideDivider
  \vspace{0.25em}
}

\newcommand{\guideSection}[1]{%
  \Needspace{6\baselineskip}%
  \par\vspace{1.1em}
  \begin{center}
    {\color{aliusline}\rule{0.14\textwidth}{0.45pt}}\hspace{0.85em}%
    {\large\bfseries #1}%
    \hspace{0.85em}{\color{aliusline}\rule{0.14\textwidth}{0.45pt}}
  \end{center}
  \vspace{-0.25em}
}

\newcommand{\guideSubsection}[1]{%
  \par\vspace{0.65em}{\bfseries #1}\par\vspace{-0.35em}
}

\newcommand{\guideQuestion}[1]{%
  \par\vspace{0.55em}{\color{aliusgreen}\bfseries #1}\par
}

\newenvironment{guideChecklist}
  {\begin{itemize}[label=\textendash]}
  {\end{itemize}}
}{\usepackage[a4paper,top=27mm,bottom=24mm,left=25mm,right=25mm]{geometry}
\usepackage[T1]{fontenc}
\usepackage[utf8]{inputenc}
\usepackage[scaled=0.96]{helvet}
\renewcommand{\familydefault}{\sfdefault}
\usepackage{array}
\usepackage{enumitem}
\usepackage{fancyhdr}
\usepackage[hidelinks]{hyperref}
\usepackage{microtype}
\usepackage{needspace}
\usepackage{tabularx}
\usepackage{xcolor}
\usepackage{xurl}

\definecolor{aliusgreen}{HTML}{13823A}
\definecolor{aliusgray}{HTML}{7A7A7A}
\definecolor{aliusline}{HTML}{9A9A9A}

\hypersetup{
  colorlinks=true,
  linkcolor=aliusgreen,
  urlcolor=aliusgreen
}

\setlength{\parindent}{0pt}
\setlength{\parskip}{0.72em}
\setlength{\tabcolsep}{0pt}
\raggedbottom
\sloppy
\urlstyle{same}

\newcommand{\ALIUSFooterTitle}{ALIUS Bulletin Guidelines (2022)}

\pagestyle{fancy}
\fancyhf{}
\fancyfoot[L]{\footnotesize\color{aliusgray}\ALIUSFooterTitle}
\fancyfoot[R]{\footnotesize\color{aliusgray}\thepage}
\renewcommand{\headrulewidth}{0pt}
\renewcommand{\footrulewidth}{0.45pt}

\setlist[itemize]{
  leftmargin=1.35em,
  itemsep=0.18em,
  topsep=0.15em,
  parsep=0em
}
\setlist[enumerate]{
  leftmargin=2.05em,
  itemsep=0.42em,
  topsep=0.2em,
  parsep=0em
}
\setlist[description]{
  leftmargin=0em,
  labelsep=0.55em,
  itemsep=0.35em,
  topsep=0.2em,
  font=\normalfont\bfseries
}

\newcommand{\guideDivider}{%
  \par\noindent{\color{aliusline}\rule{\textwidth}{0.45pt}}\par
}

\newcommand{\guideMeta}[2]{%
  {\footnotesize\bfseries\MakeUppercase{#1}}\par
  {\footnotesize #2}\par\vspace{0.65em}
}

\newcommand{\guideHeader}[5]{%
  \thispagestyle{fancy}%
  \vspace*{1mm}
  \begin{center}
    {\Large\bfseries\MakeUppercase{ALIUS Bulletin}\par}
    \vspace{0.25em}
    {\normalsize *\par}
    \vspace{0.25em}
    {\normalsize\bfseries\MakeUppercase{Guidelines For}\par}
    \vspace{0.15em}
    {\LARGE\bfseries\MakeUppercase{#1}\par}
  \end{center}
  \vspace{1.05em}
  \begin{tabularx}{\textwidth}{@{}>{\raggedright\arraybackslash}p{0.55\textwidth}!{\hspace{1.35em}{\color{aliusline}\vrule width 0.5pt}\hspace{1.35em}}>{\raggedright\arraybackslash}X@{}}
    {\large\itshape #2} &
    \guideMeta{Format}{#3}
    \guideMeta{Editorial Route}{#4}
    \guideMeta{Contact}{#5}
  \end{tabularx}
  \vspace{0.75em}
  \guideDivider
  \vspace{0.25em}
}

\newcommand{\guideSection}[1]{%
  \Needspace{6\baselineskip}%
  \par\vspace{1.1em}
  \begin{center}
    {\color{aliusline}\rule{0.14\textwidth}{0.45pt}}\hspace{0.85em}%
    {\large\bfseries #1}%
    \hspace{0.85em}{\color{aliusline}\rule{0.14\textwidth}{0.45pt}}
  \end{center}
  \vspace{-0.25em}
}

\newcommand{\guideSubsection}[1]{%
  \par\vspace{0.65em}{\bfseries #1}\par\vspace{-0.35em}
}

\newcommand{\guideQuestion}[1]{%
  \par\vspace{0.55em}{\color{aliusgreen}\bfseries #1}\par
}

\newenvironment{guideChecklist}
  {\begin{itemize}[label=\textendash]}
  {\end{itemize}}
}
\renewcommand{\ALIUSFooterTitle}{ALIUS Bulletin Guidelines for Interviews (2022)}

\begin{document}

\guideHeader
  {Interviews}
  {If you want to conduct an interview for the ALIUS Bulletin, please read the following guidelines carefully.}
  {Written scholarly interview; at least 3,500 words including questions, responses, and bibliography.}
  {Rolling publication when ready, followed by annual Bulletin integration.}
  {\href{mailto:aliusresearchgroup@gmail.com}{aliusresearchgroup@gmail.com}}

Our aim for this Bulletin is to include at least five interviews.

\guideSection{New Schedule System}

For this Bulletin, we introduce a rolling publication system. Interviews are published whenever they are ready, that is, when all stages of the publication process are fulfilled. All interviews are integrated at the end of the year in the annual Bulletin. Page numbers and citation format are updated in the final publication.

\guideSection{Length of the Interview}

Each interview is expected to be at least 3,500 words long. This includes the interviewer's questions, the interviewee's responses, and the bibliography.

\guideSection{Content of the Interview}

You are not expected to carry out the interview in a journalistic style, asking short and non-technical questions. Think of this interview almost as a dialogue. If needed, do not hesitate to ask long and detailed questions.

Bear in mind that interviewees do not necessarily have a lot of time available and might take many things for granted. Therefore, your job, as the interviewer, is to make sure that your questions summarize key concepts and facts in an understandable manner, facilitating the reading for non-specialists. When relevant, do not hesitate to write a few lines introducing key concepts and key data to contextualize the questions you are asking. That being said, if you want to ask short questions, feel free to do so: write long questions only when the intelligibility of what you are saying requires it.

Generally speaking, the interviews are intended for scholars. Therefore, you are not expected to ask general-audience questions as a science journalist would do. However, each interview should be easily understandable to any scholar. For example, an anthropologist reading an interview conducted with a neuroscientist should easily understand most of the discussion, and conversely, a neuroscientist reading an interview conducted with an anthropologist should easily understand most of the discussion.

We invite you to look at Fortier and Friedman's interview with Friston or Limanowski and Milli\`ere's interview with Metzinger in the second issue of the Bulletin to see concrete examples of how questions can be phrased. However, these two models are just illustrations of what can be done. Feel free to differ from these models: as long as you broadly follow the above recommendations, it will be perfect.

\guideSection{Modus Operandi and Requirements}

\begin{enumerate}
\item Ask one of the editors of the Bulletin whether the researcher you want to interview fits well with the Bulletin's editorial line.
\item Ask the researcher you want to interview whether they agree in principle to be interviewed. When you do so, let them know about the general concept of the Bulletin, about the other researchers interviewed in previous issues of the Bulletin, and specify that the interview is done in a written way. Describe the modus operandi, the schedule, and so on.
\item You can now write up your questions. Feel free to discuss them with one of the editors if needed, but this is not required.
\item Send your questions to the interviewee. When you do so, set up a deadline specifying when the interviewee is expected to send back their responses.
\item This step is optional: once you receive the interviewee's responses, and if time permits, you may send a second round of questions. Some of the interviewee's responses may call for further questions. It is up to you to decide whether to send a new series of questions.
\item Once you have all the responses from the interviewee, you will next have to:
  \begin{itemize}[label=\textendash,leftmargin=1.5em]
  \item Give a title to the interview. It is better to discuss this directly with the interviewee, as you will need their final agreement on the title and the content of the interview.
  \item Write a short abstract, under 200 words, and up to five keywords for referencing the interview. You will need the interviewee's final agreement on these.
  \item If the interviewee discusses books or articles in the interview but does not provide the detailed reference, insert a bibliography at the end of the interview where all books and articles mentioned by you, in your questions, or by the interviewee, in their responses, are referenced.
  \item Select two to six excerpts, depending on the length of the interview, highlighting key points made by the interviewee and nicely epitomizing their ideas. These passages must be no longer than one or two sentences. Insert the selected excerpts at the very end of the text.
  \end{itemize}
\item When the file is ready, send the whole file to the editors. The file must include the interview, a title, a bibliography, and selected excerpts, and must comply with the following requirements:
  \begin{itemize}[label=\textendash,leftmargin=1.5em]
  \item Use Times New Roman, size 12.
  \item All your questions must appear in red; all the interviewee's responses must appear in black.
  \item Citation style for the bibliography is the 6th edition of APA. If you use Zotero, this citation style can be downloaded from the \href{https://www.zotero.org/styles?q=id%3Aapa}{Zotero APA style page}.
  \item The Bulletin is written in American English rather than British English. If the text that you or the interviewee wrote is in British English, it will have to be turned into American English, for example: programme becomes program, realised becomes realized, and so on. If you or the interviewee uses British lexicon or phrasing, that is fine: you can keep it as it is. Typographic details must still follow American English. Keep in mind that in American English, abbreviations such as ``e.g.'' and ``i.e.'' are followed by a comma, for example: ``i.e., the mind is\ldots'' and not ``i.e. the mind is\ldots''
  \item Use the following quotation marks: `` ''. Do not use guillemets or single quotation marks as the main quotation marks. If there is a quotation within a quotation, use single quotation marks. For example: ``I know that the meaning he gives to `consciousness' differs from\ldots''
  \item Make sure there is no space before punctuation marks such as semicolons, colons, commas, periods, exclamation marks, or question marks.
  \item Dashes used for punctuation must be em dashes and not en dashes, and they must not be separated by spaces. For example: ``the concept of consciousness---thus defined---proves difficult.''
  \item Before sending the file to the editors, take time to carefully go through the interview to make sure there is no typo.
  \end{itemize}
\item Once they receive the text of the interview, the editors will forward it to a native English-speaking proofreader who will check grammar and linguistic correctness. Note that the proofreader's job is not to correct typos. Typos must be corrected before sending the file to the editors.
\item The editors will next send you the feedback provided by the proofreaders. You will be expected to integrate the suggested changes into the interview.
\item Once all changes have been integrated, read the whole interview again very carefully to make sure there is no typo left.
\item When the final version is ready, send it to the interviewee to get their final agreement for publication.
\item Once you have the interviewee's agreement for publication, send the final version of the interview to one of the editors. The interview will then be ready to be published.
\item If the interviewee does not agree to publish the interview as it is and wants changes to be made, make these changes and then send the revised version of the piece to the editors. They will forward it once more to the proofreaders to make sure there are no grammatical or linguistic problems. Once proofreaders have gone through the interview, proceed as described above.
\end{enumerate}

\guideSection{Editorial Team}

For any questions, feel free to ask the editors: Daniel Friedman (\href{mailto:danielarifriedman@gmail.com}{danielarifriedman@gmail.com}); Matthieu Koroma (\href{mailto:matthieu.koroma@gmail.com}{matthieu.koroma@gmail.com}); George Fejer (\href{mailto:gyorgyf@live.com}{gyorgyf@live.com}).

\end{document}
